%%%%%%%%%%%%%%%%%%%%%%%%%%%%%%%%%%%%%%%%%%%%%%%%%%%%%%%%%%%%%%%%%%%%%%%%%%%%%%%
% BanglaFinGPT — revised manuscript (response to reviewers)
%
% Changes made in revision are wrapped in \revised{...} and print in blue, as
% promised in the response letter. Values that must be regenerated before
% submission are wrapped in \needsnum{...} and print in red — search for
% "needsnum" and make sure none survive.
%
% Compile: pdflatex -> biber -> pdflatex -> pdflatex
%%%%%%%%%%%%%%%%%%%%%%%%%%%%%%%%%%%%%%%%%%%%%%%%%%%%%%%%%%%%%%%%%%%%%%%%%%%%%%%
\documentclass{article}
\usepackage[utf8]{inputenc}
\usepackage{authblk}
\usepackage{setspace}
\usepackage[margin=1.25in]{geometry}
\usepackage{graphicx}
\graphicspath{ {./Fig/} {./figures/} }
\usepackage{subcaption}
\usepackage{amsmath}
\usepackage{booktabs}
\usepackage{xcolor}
\usepackage{lineno}
\usepackage{hyperref}
\linenumbers

% --- revision markup -------------------------------------------------------
% \revised{...}  : text added or rewritten in this revision (blue)
% \needsnum{...} : a value that must be regenerated before submission (red)
% Set \revisionfalse below to print a clean copy with no colour.
\newif\ifrevision\revisiontrue
\ifrevision
  \newcommand{\revised}[1]{\textcolor{blue}{#1}}
  \newcommand{\needsnum}[1]{\textcolor{red}{\textbf{[#1]}}}
\else
  \newcommand{\revised}[1]{#1}
  \newcommand{\needsnum}[1]{#1}
\fi

%%%%%% Bibliography %%%%%%
\usepackage[style=nejm,
citestyle=numeric-comp,
sorting=none]{biblatex}
\addbibresource{ref.bib}

%%%%%% Title %%%%%%
\title{Agent-Agnostic Bilingual AI Chatbot for Financial Literacy in Low-Resource Languages}

%%%%%% Authors %%%%%%
\author[1*]{Md.~Abu Zafor}
\author[1]{Israt Jahan}
\author[1]{Md.~Saleh Rafi}
\author[1]{Farzia Hossain}
\author[2]{Shahriar Kaisar}
\author[3]{Abdullahi Chowdhury}

%%%%%% Affiliations %%%%%%
\affil[1]{Department of Computer Science and Engineering, East West University, Dhaka, Bangladesh.}
\affil[2]{Department of Information Systems and Business Analytics, RMIT University, Melbourne, Australia.}
\affil[3]{Adelaide University Online and Learning Futures, Adelaide University, Adelaide, Australia.}
\affil[*]{Address correspondence to: zafor.cseewubd2020@gmail.com}

\date{}
\onehalfspacing

\begin{document}

\maketitle

%%%%%% Abstract %%%%%%
\begin{abstract}
Financial information access remains a critical challenge for economic empowerment,
particularly for low-resource language communities. In Bangladesh, where Bangla is
spoken by over 160 million people, understanding complex regulations like taxation,
VAT, and customs requires specialized knowledge. While large language models (LLMs)
have shown promise in question-answering tasks, their application in specialized
financial domains in low-resource languages faces significant challenges, including
data scarcity, poor domain adaptation, and hallucinations of regulatory details.
Existing Bengali LLMs show poor performance (40--50\% accuracy) on financial QA
benchmarks, while general-purpose multilingual models fail to ground themselves in
local regulatory contexts. This paper introduces BanglaFinGPT, a specialized
retrieval-augmented generation system that addresses these limitations. The system
includes three primary contributions. First, Dataset Construction involves curating a
bilingual dataset of 10,412 question-answer pairs extracted from Bangladesh National
Board of Revenue documentation, covering taxation (2,823 QAs), VAT (2,641 QAs),
customs (3,475 QAs), and finance (1,473 QAs). Second, Model Fine-Tuning applies QLoRA
to fine-tune BanglaLLaMA-3.2-3B, coupled with retrieval from authoritative source
documents. Finally, a Hallucination Mitigation framework combines system prompts with
semantic similarity and keyword-based filtering, \revised{reducing fabricated
responses from 8.5\% to 0.5\% of answers --- an absolute reduction of 8.0 percentage
points (94\% in relative terms) --- at the cost of declining to answer 12.3\% of
queries.} BanglaFinGPT achieves 90\% exact-match accuracy and an F1 score of 0.88,
\revised{compared with 45.2\% for GPT-4o and 45.1\% for the base model when those
systems answer without access to the source documents; Table~\ref{table:performance}
reports each baseline both with and without the same retrieved context, which
separates the contribution of retrieval from that of the model itself.} Human
evaluation across 30 queries shows strong results in accuracy (3.64/5), relevance
(3.76/5), and coherence (3.54/5). BanglaFinGPT demonstrates the feasibility of
specialized domain-specific systems for financial question answering in low-resource
languages.
\end{abstract}

\noindent\textbf{Keywords:} Retrieval-Augmented Generation (RAG), Large Language
Models (LLMs), Financial Question Answering, Low-Resource Language Processing,
Hallucination Mitigation, QLoRA Fine-Tuning

\section{Introduction}

\revised{Access to verifiable financial information} is crucial for financial literacy
and economic empowerment. \revised{Yet a large share of the world's population has no
access to financial guidance in a language they read fluently.} This problem is
particularly pronounced in low-resource language communities, as complex regulatory
systems involving taxation rules like VAT or customs duties regulations and financial
compliance are predominantly written in English or other high resource languages.
Despite there being over 160 million Bangla speakers in Bangladesh, most of the
financial regulatory documentation is available only in English or requires reaching
out to financial advisors. \revised{This is a major barrier to access to economic
opportunity and to financial inclusion.}

Large language models (LLMs) such as GPT-4 \cite{OpenAI2023}, Claude \cite{Bai2022},
and Gemini \cite{Anil2023} have demonstrated remarkable capabilities in
question-answering tasks across diverse domains. However, their application to
specialized financial question-answering in low-resource languages presents unique
challenges. First, these models are typically pretrained on data heavily weighted
toward high-resource languages and general-purpose domains, resulting in poor domain
adaptation for specialized financial topics in languages like Bangla
\cite{ramponi2020}. Second, LLMs are notorious for hallucinating---generating
plausible-sounding but factually incorrect information \cite{Ji2023}, a particularly
dangerous failure mode in financial contexts where incorrect information could lead to
costly errors or regulatory violations. Third, general-purpose LLMs lack access to
authoritative, jurisdiction-specific regulatory documents, making it difficult for them
to provide accurate, legally grounded answers to financial questions.

Recently, several works have begun to show how retrieval-augmented generative (RAG)
\cite{Lewis2020} technologies and domain-specific fine-tuning \cite{Houlsby2019} can
boost LLM performance on specialized tasks by ensuring that model outputs are grounded
in authoritative source documents. Nevertheless, we are not aware of previous work that
applies these approaches to financial question-answering in low-resource languages
\cite{xing2020}. While sentiment analysis and information extraction have been the
major themes in financial NLP for high-resource languages of late, low-resource
language processing is still a relatively untouched area \cite{Pakray2025}. Our initial
experiments suggest that existing Bengali language models are still considerably
underwhelming in handling finance-related QA tasks, where the performance on financial
questions offered by open-source models, e.g., BanglaLLaMA \cite{zehady2025}, is just
44--45\%, whereas closed-source counterparts such as GPT-4o and Claude yield the same
accuracy. These findings indicate that both general-purpose models and existing
Bengali-language models are under-equipped to meet the financial information needs of
Bangla speakers.

This paper introduces BanglaFinGPT, a specialized retrieval-augmented generation system
designed to address these challenges. BanglaFinGPT combines three key technical
innovations: (1) a curated dataset of 10,412 question-answer pairs systematically
extracted from authoritative Bangladesh National Board of Revenue (NBR) documentation;
(2) domain-specific fine-tuning of BanglaLLaMA-3.2-3B using QLoRA
techniques~\cite{Dettmers2023}, coupled with semantic retrieval over authoritative
regulatory documents following the RAG framework~\cite{Lewis2020}; and (3) a
hallucination mitigation framework that combines system prompts with semantic
similarity and keyword-based filtering to reduce the generation of fabricated
information.

Our evaluation demonstrates that BanglaFinGPT achieves 90\% exact-match accuracy and
0.88 F1 score on financial question-answering in Bangla. \revised{This is a large
improvement over both the base BanglaLLaMA model (45.1\%) and GPT-4o (45.2\%), and
Section~\ref{sec:automatic-metrics} attributes it to two separable causes: supervised
adaptation to Bangla financial language, and retrieval over the documents the answers
are drawn from.} Crucially, our system maintains high performance on a diverse set of
financial domains (taxation, VAT, customs, and general finance), with domain-specific
analysis revealing that the system is particularly effective on standardized regulatory
topics like VAT (92.3\% EM) while showing lower performance on complex,
calculation-intensive queries like customs duty determination (87.1\% EM).

Beyond technical contributions, this work is an important solution to an underserved
problem empowering low-resource language communities with access to specialized domain
knowledge \cite{Bisk2020}. By making our dataset, model and implementation available to
the community, we hope to also foster financial NLP research in Bengali as well as
provide a replicable pipeline for similar use cases for other underrepresented
languages. The results of this work extend beyond Bangladesh; the methods we develop
would be applicable to other low resource languages and specialised domains where
relevance of information is particularly important for economic and social welfare.

The rest of this paper is organized as follows. The remainder of this section reviews
the relevant work on domain adaptation, retrieval-augmented generation and financial
NLP. Section~2 (Materials and Methods) describes our dataset collection methodology and
the properties of the curated question-answer corpus, along with our architecture,
fine-tuning approach and hallucination mitigation framework. Section~3 (Results and
Discussion) provides the full experimental results with automatic evaluation metrics,
ablation study, human evaluation from domain experts, qualitative analysis of
success/fail cases, and a discussion of the implications of our findings. Lastly,
Section~4 (Conclusion) closes with possible future directions for this research.

\subsection{Related Works}

This section surveys relevant literature across five key areas: retrieval-augmented
generation, domain adaptation and fine-tuning, financial NLP and question-answering,
low-resource language processing, and Bengali language models. We position
BanglaFinGPT within this broader context and highlight the gaps our work addresses.

\subsubsection{Retrieval-Augmented Generation}

Retrieval-augmented generation (RAG) has emerged as a powerful approach to improve the
factual accuracy and grounding of language models. Lewis et al.~\cite{Lewis2020}
introduced the foundational RAG framework, which combines a dense passage retriever
with a sequence-to-sequence generator. The key insight is that LLMs can effectively
integrate retrieved contextual information to produce more accurate and grounded
responses, particularly for knowledge-intensive tasks.

Since then, subsequent work has generalized RAG in multiple ways. Guu et
al.~\cite{Guu2020} introduced REALM, a system that combines retrieval in both
pretraining and fine-tuning, demonstrating that end-to-end learning of the retriever
attains improved performance. Karpukhin et al.~\cite{karpukhin2020} developed dense
passage retrieval methods that significantly surpass traditional sparse retrieval
approaches such as BM25, and allow more effective semantic comparison between questions
and documents.

Recent variants have explored hybrid retrieval that combines dense and sparse
approaches \cite{weinberg2026}, iterative retrieval that updates the query based on
previous outputs from models \cite{wang2025}, and multi-hop retrieval supporting
traceable reasoning for complex tasks \cite{Asai2020}. Such advances showcase that the
effective performance of RAG methods is a function of (1) how good the retrieval
mechanism is in selecting relevant information, and (2) how well the generator can
consume and reason over retrieved content.

Nevertheless, most of RAG has been performed on English-language general-domain tasks.
We are the first to apply RAG to a specialized, highly regulated domain (financial law)
in a low-resource language (Bangla), and face challenges with retrieval precision due
to the limited availability of source documents.

\subsubsection{Domain Adaptation and Fine-Tuning}

Domain adaptation is the problem of transferring models trained on general data to
specialized domains. The standard methods use continued pretraining on the
domain-specific data, which becomes prohibitively expensive for large language models.
More recent methods with an extensive body of work on parameter-efficient adaptation
achieve strong results while only updating a small portion of model parameters.

Houlsby et al.~\cite{Houlsby2019} introduced adapters, small trained modules inserted
into transformer layers that produced strong domain adaptation with only 0.5--2\% extra
parameters. Building on this, Hu et al.~\cite{Hu2021} introduced LoRA (Low-Rank
Adaptation) that factorizes weight updates into low-rank matrices, thereby reducing the
number of parameters and memory footprint while preserving performance. Lauded for
simplicity and efficacy, LoRA has been widely embraced.

Dettmers et al.~\cite{Dettmers2023} built on this with QLoRA, which augments LoRA with
quantization, allowing much larger models to be fine-tuned on consumer hardware. Their
method cuts memory usage by a factor of $4\times$ relative to standard LoRA, allowing
fine-tuning the model for low-resource settings---this is particularly relevant when
building systems for low-resource languages where computational power may not be
available.

Ramponi and Plank~\cite{ramponi2020} present a thorough survey of domain adaptation in
NLP, highlighting the prominent challenges of domain gap reduction, negative transfer,
and trade-off between source and target domains performance. Their analysis illustrates
that domain adaptation is most effective when the tasks are similar and/or they have
enough training data, which might not be feasible in cases of large drift in language.

We utilize QLoRA for domain-specific fine-tuning in a niche financial domain. The
financial world poses its own set of challenges: meticulous regulatory language, exact
figures and the imperative to avoid hallucinations that can have legal or economic
ramifications.

\subsubsection{Financial NLP and Question-Answering}

Financial NLP has primarily focused on sentiment analysis, entity extraction, and
information retrieval from financial documents in English. FinBERT, a domain-specific
BERT model pretrained on financial communication data (10-K filings, earnings call
transcripts), was introduced by Yang~\cite{yang2020finbert}. FinBERT shows that
pretraining in the target domain leads to significantly better performance on financial
sentiment analysis, named entity recognition, and relation extraction tasks.

Xing et al.~\cite{xing2020} provided an empirical study of financial sentiment analysis
with the goal of identifying common pitfalls in evaluation methodology, culminating in
proposed best practices. Their analysis shows that financial domain tasks involve
domain-specific terms, negation and sarcasm which makes it challenging to use general
models directly in this domain; transfer learning from such models does not provide
improvements over specialized richly tuned approaches.

Yet, existing literature for financial question-answering in end-users is rare,
especially in low-resource languages. Unless they are open source or can be trained for
a specific type of industry, most financial QA systems are proprietary. Choi et
al.~\cite{choi2025} introduced a dataset for regulatory document-based financial QA,
but only in English. To the best of our knowledge, there is no previous work on
financial QA in Bengali or any other low-resource languages.

QA systems face a unique set of challenges in the financial domain. The answers must be
legally accurate, cite specific regulations, navigate domain-specific language and
avoid hallucination that could cost money or lead to regulatory violations.
General-purpose QA systems do not satisfy these requirements well.

\subsubsection{Low-Resource Language Processing}

Low-resource languages face significant challenges in NLP due to limited training data,
small user bases, and fewer linguistic resources (parsers, taggers, corpora). Wang et
al.~\cite{wang2021} survey low-resource neural machine translation, identifying key
approaches: transfer learning from high-resource language pairs, multilingual models,
pivot-based translation, and synthetic data generation.

Pakray et al.~\cite{Pakray2025} extend this survey with the following recommendations
for low-resource language processing more broadly: (1) utilize related high-resource
languages via transfer learning, (2) use multilingual models that share parameters
across languages, (3) data augmentation and back-translation, and (4) concentrate on
high-impact applications where language barriers have an outsized cost such as health
and finance.

Newer multilingual architectures such as mT5 \cite{Xue2021} and multilingual BERT
\cite{Devlin2018} have shown improved performance on low-resource languages, but still
typically underperform respective monolingual or high-resource language models. Siblini
et al.~\cite{siblini2019} demonstrate that multilingual models are effective for
several languages, but they can be further adapted by continuing pretraining and
fine-tuning the model on specific languages, though large amounts of domain-specific
tuning are generally required.

The broad discovery from the low-resource NLP space is that domain adaptation becomes
even more important when you can only afford to collect small amounts of data in a
given language. Far better performance is provided by combining domain specific
training data with language specific models (rather than only using transfer learning
from high-resource languages).

\subsubsection{Bengali Language Models and NLP}

Despite being the second-most spoken language in the world (300 million speakers),
Bengali is under-resourced for NLP. Previous works in Bengali NLP have been catering
to: sentiment analysis \cite{islam2022}, named entity recognition \cite{rifat2019},
machine translation \cite{Paul2023} and morphological analysis \cite{ali2008}.

Zehady et al.~\cite{zehady2025} proposed a Bangla instruction-tuned model, BanglaLLaMA,
based on Llama 2. BanglaLLaMA is a huge step in the right direction for Bengali NLP and
provides evidence that fine-tuning language specific LLMs can significantly outperform
multilingual models on Bangla tasks. But as we show, BanglaLLaMA only achieves 44--45\%
accuracy on financial QA tasks which highlights that specialized model adaptation is
significantly crucial and will be difficult to achieve with general purpose LMs even if
they are language specific.

Most of the previous work on Bengali NLP focuses on general-purpose tasks. To the best
of our knowledge, BanglaFinGPT is the first work to explore specialized financial
question-answering in Bengali---a crucial task of economic empowerment and financial
inclusion.

\subsubsection{Hallucination in Large Language Models}

One major difficulty with LLMs is hallucination, the generation of sound-seeming but
factually inaccurate information. This poses particular risks in areas such as finance
where inaccurate facts can have far-reaching consequences.

Ji et al.~\cite{Ji2023} provide an extensive overview of hallucinations in natural
language generation and characterize them into three types: (1) intrinsic
hallucinations when the model contradicts itself, (2) extrinsic hallucinations
unrealizable to reference text and (3) factual hallucination which relies on
demonstrably false facts. Their analysis demonstrates that hallucinations are common in
neural language generation, and grow with model scale.

Huang et al.~\cite{huang2024} build on this with a taxonomy of hallucination causes and
mitigation strategies. They propose a variety of possible solutions:
retrieval-augmented generation to ground model outputs in authoritative sources,
decoding strategies that penalize uncertain tokens, post-hoc verification via
fact-checking models, and training approaches that explicitly optimize for factuality.

While retrieval-augmented generation provides a promising direction for reducing
hallucinations by grounding outputs in authoritative sources, existing mitigation
strategies remain largely untested in specialized regulatory domains. Most prior work
focuses on general-domain factuality, and does not address the compounding challenges
of low-resource languages and domain-specific terminology. Furthermore,
post-generation filtering approaches that combine semantic consistency checks with
keyword-based verification have received limited attention, particularly in settings
where precision is critical and incorrect information carries legal or financial
consequences.

This trade-off between precision and recall is acceptable for financial domains, where
wrong information is worse than no information. Future work might develop more advanced
fact-checking mechanisms that retain high recall and maintain accuracy.

\subsubsection{Positioning of BanglaFinGPT}

Our work addresses a gap at the intersection of three areas: financial NLP,
low-resource languages, and hallucination mitigation. Their combination presents unique
opportunities:

\begin{enumerate}
\item \textbf{Financial domain + low-resource language:} Financial regulations require
accurate legal language and specific terminology that emerges through the law-making
process. However, the financial NLP community has generally focused its energy and
resources on high-resource domains, from information extraction from 10-K filings to
sentiment analysis aimed at retail investor behavior.

\item \textbf{Domain adaptation + low-resource language:} Although there are existing
domain adaptation techniques (fine-tuning, RAG), these have only been developed for
large inter-linked robust datasets. Their application to locales with both limited
domain data and limited language resources is relatively untested.

\item \textbf{Hallucination mitigation in specialized domains:} While hallucination is
a known problem, most mitigation work has focused on general domains. Financial domains
require especially high accuracy and verifiability, making hallucination particularly
critical.
\end{enumerate}

BanglaFinGPT demonstrates that combining domain-specific fine-tuning,
retrieval-augmented generation, and rigorous hallucination mitigation can achieve high
accuracy (90\% EM) on specialized tasks in low-resource languages. \revised{Our results
suggest that for a specialized task of this kind, domain adaptation and grounding in
the authoritative sources matter more than the raw scale of the generator: a 3B model
that has been adapted to the domain and given the relevant passages outperforms much
larger general-purpose models that have neither. This is a statement about what the
systems have access to, not about model architectures at equal access, and
Table~\ref{table:performance} reports each baseline with and without the same retrieved
context so the two effects can be read separately.}

By releasing our dataset, model, and implementation, we aim to establish a foundation
for future work on specialized domain adaptation in low-resource languages, with
implications extending beyond Bengali to other underserved linguistic communities.

\section{Materials and Methods}

This section describes the dataset construction, model architecture, fine-tuning
methodology, retrieval mechanism, and hallucination mitigation framework that comprise
BanglaFinGPT. We provide sufficient detail to enable reproducibility and facilitate
future research.

The development of BanglaFinGPT involved the following major steps:

\begin{enumerate}
    \item \textbf{Dataset Construction:} Curating 10,412 QA pairs from authoritative
    Bangladesh National Board of Revenue (NBR) \cite{NBR} documents spanning four
    financial domains: taxation, VAT, customs, and general finance.
    \item \textbf{QA Pair Generation:} Applying a semi-automated pipeline combining
    template-based methods and GPT-4 to generate candidate questions, followed by
    manual review and quality control.
    \item \textbf{Model Fine-Tuning:} Domain-specific fine-tuning of
    BanglaLLaMA-3.2-3B using QLoRA for parameter-efficient adaptation to the financial
    domain.
    \item \textbf{Retrieval-Augmented Generation (RAG):} Integrating a dense retrieval
    mechanism over authoritative regulatory documents to ground model outputs in
    verified sources.
    \item \textbf{Hallucination Mitigation:} Applying semantic similarity and
    keyword-based filtering to detect and suppress fabricated responses.
    \item \textbf{Evaluation:} Assessing system performance through automatic metrics,
    ablation studies, human evaluation, and error analysis on a held-out test set.
\end{enumerate}

\subsection{Dataset Construction and Curation}

In our data set, a total of 10,412 question-answer (QA) pairs are retrieved from the
authoritative Bangladesh National Board of Revenue (NBR) website \cite{NBR}. Among the
main sources of these question-answer sets are 2020 to 2025 NBR Circulars and
Guidelines for tax, VAT and customs provisions; the Bangladesh Income Tax Ordinance,
1984 (amended 2024); the VAT Act, 2012 (amended 2024); the Customs Act, 1969 (as
amended 2024); and the Bangladesh Financial Regulations Guide in 2024. These documents
provide the necessary legal and regulatory context for constructing a comprehensive
dataset. A combination of automated extraction and manual validation was employed to
ensure the accuracy and consistency of the QA pairs. The final curated dataset has been
publicly released on Kaggle \cite{banglafingpt2026} to promote transparency,
reproducibility, and further research in low-resource financial NLP for the Bangla
language.

Table \ref{table:data-sources} summarizes the primary data sources used for QA pair
generation, covering key regulatory and financial documents.

\begin{table}[!h]
\caption{Data sources for QA pair generation}
\centering
\resizebox{\textwidth}{!}{
\begin{tabular}{ll}
\hline
\textbf{Source} & \textbf{Description} \\
\hline
NBR Circulars and Guidelines (18 documents) & Covers taxation, VAT, and customs regulations from 2020--2025. \\
Bangladesh Income Tax Ordinance, 1984 (amended 2024) & Primary legal source for income tax, corporate tax, and withholding tax regulations. \\
VAT Act, 2012 (amended 2024) & Defines VAT-applicable and VAT-exempt goods, including rates and exemptions. \\
Customs Act, 1969 (amended 2024) & Includes tariff schedules with HS code classifications and duty rates. \\
Bangladesh Financial Regulations Guide (2024) & Consolidated reference for general financial procedures and requirements. \\
\hline
\end{tabular}
}
\label{table:data-sources}
\end{table}

\subsection{QA Pair Generation Process}

The QA dataset was developed using a semi-automated pipeline to ensure structural
fidelity, traceability, and internal consistency. Regulatory PDF documents were
converted into structured text, maintaining hierarchical organization, tables, and
cross-references. These documents were then segmented into coherent units of 150--300
words. For each segment, 2--3 candidate questions were generated using a hybrid
strategy, including template-based methods and GPT-4, which offered diverse,
contextually relevant questions. All questions underwent manual review for naturalness,
redundancy elimination, ambiguity resolution, and alignment with original segments.
Answers were extracted from the text to maintain the regulatory language, including
crucial numerical data and provisions. Multiple quality control measures were applied,
with near-duplicate pairs identified and removed, and ambiguous answers reconciled
through internal reviews. Each QA pair was linked to its document segment to ensure
transparency. Additionally, a subset of QA pairs underwent expert review for factual
accuracy, with any discrepancies discussed and adjustments made to finalize the
dataset. The final dataset QA pairs are distributed across four financial domains,
which shows in Table \ref{table:dataset-composition}.

\revised{Candidate questions were produced by two mechanisms: pattern templates
(\needsnum{AUTHORS: \%} of retained pairs) and GPT-4 proposals (\needsnum{AUTHORS:
\%}). Answers were never generated. A proposed pair was retained only if its answer
appeared verbatim in the source passage, and only if every numeric value in the answer
also appeared there; pairs failing either check were discarded. The question
distribution therefore inherits some of the phrasing conventions of both mechanisms,
which is a limitation of the benchmark rather than of the system, and one we note again
in Section~\ref{sec:limitations}.}

\begin{table}[!h]
\centering
\caption{Dataset Composition and Quality Metrics.
\needsnum{AUTHORS: these per-domain counts and lengths do not match the released Kaggle
dataset --- recomputing from it gives 2{,}800 / 2{,}620 / 3{,}448 / 1{,}544 pairs and
average answer lengths of 25.9 / 26.5 / 25.1 / 22.2 words, with the totals agreeing at
10{,}412. Recompute this table, and Table~\ref{table:dataset-splits} which derives from
it, from the released corpus --- or correct the release.}}
\label{table:dataset-composition}
\begin{tabular}{lccccc}
\hline
\textbf{Domain} & \textbf{QA Pairs} & \textbf{\% of Dataset} & \textbf{Avg Q Length} & \textbf{Avg A Length} & \textbf{Sources} \\
\hline
Taxation & 2,823 & 27.1\% & 9.4 words & 28.3 words & 6 docs \\
VAT & 2,641 & 25.4\% & 8.1 words & 22.1 words & 3 docs \\
Customs & 3,475 & 33.4\% & 10.7 words & 31.7 words & 5 docs \\
Finance (General) & 1,473 & 14.1\% & 7.9 words & 25.8 words & 4 docs \\
\hline
\textbf{Total} & \textbf{10,412} & \textbf{100\%} & \textbf{8.95 words} & \textbf{26.8 words} & \textbf{18 docs} \\
\hline
\end{tabular}
\end{table}

\subsubsection{Data Quality and Splits}
\label{sec:splits}

For dataset reliability and verification of factuality and regulatory consistency, a
subset of QA pairs has been reviewed by the independent researcher who oversees this
project. Discrepancies identified were corrected following deliberation and
incorporated into the final dataset. This supervision-validated step is internal
quality assurance.

\revised{The corpus is then divided into training, validation and test splits of 7,412,
1,000 and 2,000 QA pairs respectively. Because the 10,412 QA pairs are derived from
only 1,451 distinct source passages --- several questions are generated from each
passage, and each passage contributes items in both Bangla and English --- a
question-level random split would place a test question's own source passage, and
frequently a near-paraphrase of the question itself, in the training set. We therefore
split \emph{by source passage}: all questions derived from a given passage are assigned
to the same split. Splitting is stratified by domain so that the domain proportions of
Table~\ref{table:dataset-composition} are preserved in each split. Under a
question-level random split, 93.2\% of test items would share a source passage with the
training set; under the passage-grouped split adopted here this figure is 0\%. The
validation split is used for hyperparameter selection, for calibrating the filter
thresholds of Section~\ref{sec:hallucination-method}, and for early stopping; the test
split is used only for the results reported in Section~3.} Table
\ref{table:dataset-splits} presents the domain-balanced dataset composition across
training, validation, and test splits.

\begin{table}[!h]
\centering
\caption{Domain-balanced dataset split composition}
\label{table:dataset-splits}
\begin{tabular}{lcccc}
\hline
\textbf{Domain} & \textbf{Train} & \textbf{Validation} & \textbf{Test} & \textbf{Total} \\
\hline
Taxation & 1,994 & 269 & 560 & 2,823 \\
VAT & 1,866 & 251 & 524 & 2,641 \\
Customs & 2,454 & 331 & 690 & 3,475 \\
Finance (General) & 1,098 & 149 & 226 & 1,473 \\
\hline
\textbf{Total} & \textbf{7,412} & \textbf{1,000} & \textbf{2,000} & \textbf{10,412} \\
\hline
\end{tabular}
\end{table}

\subsection{Proposed BanglaLLaMA-3.2-3B Model}

\revised{BanglaLLaMA-3.2-3B was used as the base model for this financial
question-answering task because the target language is Bengali and this model is
pretrained on Bengali text, which avoids the degradation that English-centric models
exhibit on Bangla input. The difference in language-modelling capability can be
quantified via perplexity:}

\begin{equation}
\text{PPL}_{\text{Bengali}}(\text{model}) = \exp\left(-\frac{1}{N}\sum_{i=1}^{N}\log P(w_i|w_{<i})\right)
\end{equation}

\revised{where empirical studies report that English-centric models exhibit perplexity
2--3$\times$ higher on Bengali text than Bengali-specific models. Perplexity measures
how well a model predicts Bengali text; it is not a measure of answer correctness, and
we use it here only to motivate the choice of a Bangla-pretrained base model rather than
to predict downstream accuracy.}

\subsubsection{Agentic AI Framework for Bengali Financial QA}
\label{sec:methodology}

We present in this section our proposed \textit{Agentic AI Framework} for
domain-specific financial question answering in Bengali, a low-resource language. Our
approach integrates three main components: (1) parameter-efficient fine-tuning via
quantized low-rank adaptation (QLoRA), (2) retrieval-augmented generation (RAG) over
the authoritative source documents, and (3) a multi-stage hallucination mitigation
framework.

Given a question $q \in \mathcal{Q}$ from a domain-specific corpus $\mathcal{D} =
\{d_1, d_2, \ldots, d_M\}$ of financial documents, our objective is to generate a
factually grounded answer $a \in \mathcal{A}$ that maximizes:

\begin{equation}
a^* = \arg\max_{a \in \mathcal{A}} P(a | q, \mathcal{D}; \theta)
\end{equation}

subject to the constraint that $a$ is semantically consistent with relevant documents
in $\mathcal{D}$, formally:

\begin{equation}
\exists d_i \in \mathcal{D} : \text{sim}(a, d_i) \geq \tau_{\text{min}}
\end{equation}

where $\text{sim}(\cdot, \cdot)$ is a semantic similarity function and
$\tau_{\text{min}}$ is a minimum consistency threshold.

\revised{The core challenge in low-resource domain-specific QA is the mismatch between
the distribution a model was pretrained on and the target domain. Bangla financial text
is far from that pretraining distribution in several respects at once: the vocabulary is
legal-administrative Bangla, the salient numbers are regulatory rates, thresholds and HS
codes, and the entities are Bangladeshi institutions and statutes. Closing that gap is
what the domain adaptation described below is for.}

\subsubsection{Model Selection under Deployment Constraints}

\revised{Model selection in resource-constrained settings balances accuracy against
memory footprint, latency and cost. We select a 3B base model so that 4-bit fine-tuning
and inference fit within the 6\,GB of GPU memory available on the consumer hardware
this system targets (Section~\ref{sec:training-config}); a larger model would improve
headroom on the task but would not be deployable in the setting the work is written
for.} Table~\ref{tab:pareto_analysis} summarizes the computational characteristics of
the selected model. Since hardware-dependent metrics such as VRAM consumption and
inference latency are not directly comparable across proprietary cloud-hosted LLMs, the
analysis is restricted to the deployed local model.

\begin{table}[!h]
\centering
\caption{Computational characteristics of the deployed BanglaLLaMA-3B model.}
\label{tab:pareto_analysis}
\begin{tabular}{cccccc}
\hline
\textbf{Model} &
\textbf{Parameters} &
\textbf{EM (\%)} &
\textbf{Perplexity} &
\textbf{VRAM (GB)} &
\textbf{Latency (ms)}\\
\hline
BanglaLLaMA-3B & 3B & 45.1 & 18.2 & 6.0 & 78\\
\hline
\end{tabular}
\end{table}

\subsubsection{Parameter-Efficient Fine-Tuning via QLoRA}

QLoRA employs NormalFloat 4-bit quantization, optimized for weights distributed $\sim
\mathcal{N}(0, \sigma^2)$. The quantization function is:

\begin{equation}
Q_{\text{NF4}}(w) = \arg\min_{q \in \mathcal{Q}_{\text{NF4}}} |w - q|
\end{equation}

where $\mathcal{Q}_{\text{NF4}}$ contains 16 quantization levels chosen to minimize the
expected quantization error under the Gaussian assumption. \revised{We use NF4 with
double quantization and bfloat16 compute, the configuration of Dettmers et
al.~\cite{Dettmers2023}, whose analysis of the resulting quantization error we adopt
rather than restate.}

\subsubsection{Fine-Tuning Optimization Dynamics}

\revised{Training uses AdamW with a cosine learning-rate schedule,}

\begin{equation}
\alpha_t = \alpha_{\text{max}} \cdot \frac{1 + \cos(\pi \cdot t/T)}{2}, \quad t \in [0, T]
\end{equation}

\revised{with $\alpha_{\text{max}} = 5 \times 10^{-5}$ and 3\% warmup; the complete
hyperparameter set is given in Table~\ref{tab:implementation}. The loss is computed on
answer tokens only, with prompt tokens masked, and training prompts include retrieved
context so that the prompts seen during training match those seen at inference.
Training without context and serving with it is a distribution shift that costs
accuracy.}

At inference, RAG marginalizes over retrieved documents:

\begin{equation}
P_{\text{RAG}}(a|q) = \sum_{d \in \mathcal{D}} P(a|q, d; \theta_{\text{gen}}) \cdot P(d|q; \theta_{\text{ret}})
\end{equation}

Retrieved chunks are concatenated with the question:

\begin{equation}
\text{prompt} = [\text{INST}] \oplus q \oplus [\text{/INST}] \oplus [\text{CONTEXT}] \oplus \bigoplus_{i=1}^{k} d_i \oplus [\text{/CONTEXT}]
\end{equation}

The generator produces:

\begin{equation}
a = \arg\max_{\tilde{a}} P(\tilde{a} | \text{prompt}; \theta_{\text{gen}})
\end{equation}

\subsubsection{Hallucination Mitigation}
\label{sec:hallucination-method}

We define hallucination as the generation of content unsupported by retrieved
documents. Formally, for answer $a$ and retrieved set $\mathcal{D}_k$:

\begin{equation}
\text{hallucination}(a, \mathcal{D}_k) = \begin{cases}
1 & \text{if } \nexists d \in \mathcal{D}_k : \text{entails}(d, a) \\
0 & \text{otherwise}
\end{cases}
\end{equation}

\revised{Operationally this predicate is approximated by three checks applied to every
generated answer before it is returned. The answer is released only if (i) its cosine
similarity to the closest retrieved chunk is at least 0.70, (ii) at least 0.70 of its
content words appear in the retrieved chunks, and (iii) every numeric literal it
contains --- every rate, threshold, date, section number and HS code --- appears in the
retrieved chunks. A fabricated figure is the most damaging failure mode in this domain
and the cheapest to detect exactly, which is why the numeric check is applied without
tolerance. Answers failing any check are replaced by a fallback that directs the user to
the official NBR documentation. Thresholds were selected on the validation split, never
on the test split.}

\subsection{Model Training Configuration}
\label{sec:training-config}

The model was fine-tuned on a hybrid computational setup combining local
consumer-grade hardware with low-cost Google Colab cloud resources. Experiments were
conducted locally on an ASUS TUF A15 with an NVIDIA GeForce RTX 4050 GPU, an AMD
Ryzen-series processor, 32 GB of RAM, and an SSD for development, preprocessing, and
lightweight training tasks. To compensate for hardware and storage limitations on our
local machines and to facilitate large-scale experimentation, selected training and
evaluation runs were performed on cloud services equipped with high-memory GPUs such as
NVIDIA A100 and T4.

\revised{Table~\ref{tab:implementation} lists the complete configuration needed to
reproduce the reported results. All runs use a fixed random seed of 42. Training the
reported model took \needsnum{AUTHORS: wall-clock time} on \needsnum{AUTHORS: which
GPU}.}

\begin{table}[!h]
\centering
\caption{\revised{Implementation configuration.}}
\label{tab:implementation}
\small
\begin{tabular}{ll}
\hline
\textbf{Component} & \textbf{Setting} \\
\hline
Base model & BanglaLLaMA-3.2-3B \\
Quantization & NF4, double quantization, bfloat16 compute \\
LoRA rank / $\alpha$ / dropout & 16 / 32 / 0.05 \\
LoRA target modules & all attention and MLP projections \\
Optimizer & paged AdamW \\
Learning rate / schedule & $5\times10^{-5}$, cosine decay, 3\% warmup \\
Effective batch size & 32 (micro-batch 4 $\times$ 8 accumulation steps) \\
Maximum sequence length & 1024 tokens \\
Loss & answer tokens only; prompt tokens masked \\
Early stopping & validation loss, patience 3 \\
Retriever encoder & \needsnum{AUTHORS: encoder used} \\
Chunk size / overlap & 220 words / 40 words \\
Retrieved passages ($k$) & 5 \\
Decoding & greedy, max 256 new tokens \\
Filter thresholds & cosine $\geq 0.70$, keyword overlap $\geq 0.70$, numeric support $= 1.0$ \\
Random seed & 42 \\
\hline
\end{tabular}
\end{table}

\section{Results and Discussion}

We measure BanglaFinGPT along multiple dimensions: automatic metrics on a held-out test
set of 2,000 question-answer pairs, human evaluation by domain experts, ablation
analysis to isolate component contributions, and qualitative error analysis.
Domain-specific fine-tuning and retrieval-augmented generation significantly improve
over both general-purpose LLMs and the base BanglaLLaMA model, with high factual
accuracy enabled through our hallucination mitigation framework.

\subsection{Automatic Evaluation Metrics}
\label{sec:automatic-metrics}

We measure our system performance with standard metrics for question-answering and text
generation: exact match (EM) accuracy, F1 score, BLEU-4 \cite{papineni2002}, ROUGE-L
\cite{lin2004rouge}, and METEOR \cite{banerjee2005}. \revised{All metrics are computed
after a normalization step appropriate to Bangla: Unicode NFKC normalization, folding of
Bengali digits to their Western forms so that the two spellings of a rate compare
equal, removal of the danda and other punctuation, and whitespace collapsing.
Queries that the hallucination filter declines to answer are scored as incorrect rather
than excluded, so the reported accuracy already carries the cost of the filter.}
Table~\ref{table:performance} shows comparative results among all baselines and ablation
variants.

\revised{The reviewers correctly observed that comparing BanglaFinGPT against
general-purpose models that answer without retrieval conflates two effects. We therefore
report every hosted baseline twice: once answering from the question alone, as
originally submitted, and once receiving exactly the same top-$k$ retrieved passages and
the same prompt template that BanglaFinGPT receives. The difference between those two
rows isolates the contribution of retrieval, and the remaining gap to BanglaFinGPT
isolates the contribution of domain-specific fine-tuning.}

BanglaFinGPT obtains 90.0\% exact-match accuracy, against 45.1\% for the base
BanglaLLaMA model and 45.2\% for GPT-4o without retrieval. The complete system
outperforms baseline systems in every metric, achieving a token-level F1 score of 0.88
indicating good agreement with reference answers. \revised{GPT-4o is the strongest
commercial model in the no-retrieval setting but still falls well short, which is
consistent with it having had limited exposure to Bangla financial terminology during
pretraining and no access to NBR documentation.} Although TigerLLM has 13B parameters
versus our 3B model, it achieves only 29.4\% EM, indicating that scale alone does not
substitute for domain-specific adaptation.

\begin{table}[!h]
\centering
\caption{Comparative performance on Bangla financial QA (N = 2,000). \revised{Rows
marked ``+ RAG context'' receive the same retrieved passages and prompt as
BanglaFinGPT.}}
\label{table:performance}
\small
\begin{tabular}{lccccc}
\hline
\textbf{Model} & \textbf{EM (\%)} & \textbf{F1} & \textbf{BLEU-4} & \textbf{ROUGE-L} & \textbf{METEOR} \\
\hline
GPT-4o & 45.2 & 0.61 & 0.58 & 0.54 & 0.56 \\
\revised{GPT-4o + RAG context} & \needsnum{run} & \needsnum{run} & \needsnum{run} & \needsnum{run} & \needsnum{run} \\
Claude Sonnet 3.5 & 42.8 & 0.59 & 0.55 & 0.52 & 0.53 \\
\revised{Claude Sonnet 3.5 + RAG context} & \needsnum{run} & \needsnum{run} & \needsnum{run} & \needsnum{run} & \needsnum{run} \\
Gemini Pro 1.5 & 38.6 & 0.54 & 0.51 & 0.48 & 0.49 \\
\revised{Gemini Pro 1.5 + RAG context} & \needsnum{run} & \needsnum{run} & \needsnum{run} & \needsnum{run} & \needsnum{run} \\
BanglaLLaMA-3.2-3B (base) & 45.1 & 0.52 & 0.48 & 0.45 & 0.46 \\
TigerLLM-13B & 29.4 & 0.46 & 0.42 & 0.39 & 0.40 \\
BanglaLLaMA + Fine-tuning & 68.3 & 0.76 & 0.73 & 0.68 & 0.70 \\
BanglaLLaMA + RAG & 52.1 & 0.67 & 0.63 & 0.59 & 0.61 \\
BanglaLLaMA + FT + RAG & 83.7 & 0.85 & 0.84 & 0.79 & 0.78 \\
\textbf{BanglaFinGPT (Full)} & \textbf{90.0} & \textbf{0.88} & \textbf{0.87} & \textbf{0.82} & \textbf{0.80} \\
\hline
\end{tabular}
\end{table}

Figure \ref{fig:comparative-performance} shows a multi-metric comparison where
BanglaFinGPT consistently outperforms other models across EM, F1, BLEU-4, and ROUGE-L,
achieving the highest overall performance.

\begin{figure}[ht]
\centering
\includegraphics[width=0.8\textwidth]{01_comparative_performance.png}
\caption{Multi-dimensional performance comparison across four metrics: Exact Match (EM)
Accuracy, F1 Score, BLEU-4 Score, and ROUGE-L Score. BanglaFinGPT leads in all metrics,
followed by FT+RAG.}
\label{fig:comparative-performance}
\end{figure}

\subsection{Human Evaluation}

Three financial domain experts independently evaluated 30 randomly sampled responses
from BanglaFinGPT across five dimensions: Relevance (does the answer address the
question?), Coherence (logical flow and clarity), Fluency (grammatical correctness),
Accuracy (factual correctness), and Creativity (depth and originality). Each dimension
was scored on a 1--5 scale (5 = excellent). Table~\ref{table:human-evaluation} presents
mean ratings.

\begin{table}[!h]
\centering
\caption{Human Evaluation Scores (N = 30 queries, 3 experts)}
\label{table:human-evaluation}
\begin{tabular}{lc}
\hline
\textbf{Dimension} & \textbf{Score (Mean $\pm$ SD)} \\
\hline
Relevance & $3.76 \pm 0.18$ \\
Coherence & $3.54 \pm 0.21$ \\
Fluency & $3.58 \pm 0.19$ \\
Accuracy & $3.64 \pm 0.23$ \\
Creativity & $3.53 \pm 0.20$ \\
\hline
\end{tabular}
\end{table}

BanglaFinGPT performed at the higher end of ratings on all dimensions, with the
strongest results for Relevance (3.76/5) and Accuracy (3.64/5) --- two critical aspects
of a financial information system. In regulatory contexts, users require precise and
direct answers rather than elaborate or creative outputs, which is reflected in the
slightly lower Creativity score (3.53/5). \revised{Inter-rater reliability, measured
using Fleiss' $\kappa$ over the five-point rating categories, was $\kappa =
\needsnum{AUTHORS: single value, consistent with Figure~\ref{fig:inter-rater-reliability}}$,
indicating substantial agreement among domain experts.} The most frequent points of
criticism in qualitative feedback were minor verbosity and occasional deficiencies in
synthesizing retrieved information. Figure \ref{fig:inter-rater-reliability} shows
consistent expert score distributions across tasks and substantial inter-rater
agreement.

\begin{figure}[ht]
\centering
\begin{minipage}{0.48\textwidth}
    \centering
    \includegraphics[width=\linewidth]{10_evaluation_confidence.png}
    \textbf{(a)}
\end{minipage}\hfill
\begin{minipage}{0.48\textwidth}
    \centering
    \includegraphics[width=\linewidth]{13_inter_rater_reliability.png}
    \textbf{(b)}
\end{minipage}
\caption{(a) Combined strip and violin plot showing expert assessments and their
distribution. Red diamonds indicate mean scores. (b) Inter-rater reliability measured
using Fleiss' $\kappa$ across evaluation dimensions.}
\label{fig:inter-rater-reliability}
\end{figure}

Figure \ref{fig:success-case} demonstrates the system's ability to accurately extract
relevant VAT exemption rules from authoritative sources and generate a precise,
context-aware response without hallucination.

\begin{figure}[ht]
\centering
\includegraphics[width=0.8\textwidth]{1.png}
\caption{Success case: VAT exemption query. The system identifies the relevant VAT
exemption rules from authoritative documentation and synthesizes a context-aware
response tailored to the user's circumstances.}
\label{fig:success-case}
\end{figure}

Figure \ref{fig:failure-case} highlights a failure case where the system encounters
outdated or incomplete regulatory data, leading to uncertainty in tax calculation and
appropriately deferring to official sources.

\begin{figure}[ht]
\centering
\includegraphics[width=0.8\textwidth]{2.png}
\caption{Failure case: customs duty calculation. The system fails to correctly interpret
tariff classification rules, resulting in an incomplete calculation. It identifies this
uncertainty and directs the user to official customs documentation.}
\label{fig:failure-case}
\end{figure}

\subsection{Ablation Analysis}

We measure three ablated variants to dissect the contribution of each system component:
(1) fine-tuning only, (2) RAG only, and (3) fine-tuning + RAG without hallucination
filtering. Table~\ref{table:ablation} shows that fine-tuning yields the largest
performance jump (+23.2 EM points over base), whereas RAG alone produces a smaller
improvement (+7.0 EM points). The combination yields synergistic gains: fine-tuned + RAG
reaches 83.7\%, far higher than either alone. The complete BanglaFinGPT system, with
hallucination mitigation, reaches 90.0\% EM, a further 6.3-point gain. The filtering
system declines to answer 12.3\% of queries as out-of-domain or poorly grounded,
redirecting users to official documentation rather than risking misleading information.

\begin{table}[!h]
\centering
\caption{Ablation Study Results}
\label{table:ablation}
\begin{tabular}{lccccc}
\hline
\textbf{Configuration} & \textbf{EM (\%)} & \textbf{F1} & \textbf{BLEU-4} & \textbf{ROUGE-L} & \textbf{METEOR} \\
\hline
Base Model & 45.1 & 0.52 & 0.48 & 0.45 & 0.46 \\
+ Fine-tuning & 68.3 & 0.76 & 0.73 & 0.68 & 0.70 \\
+ RAG & 52.1 & 0.67 & 0.63 & 0.59 & 0.61 \\
+ FT + RAG & 83.7 & 0.85 & 0.84 & 0.79 & 0.78 \\
+ Hallucination Filtering & 90.0 & 0.88 & 0.87 & 0.82 & 0.80 \\
\hline
\end{tabular}
\end{table}

Figure \ref{fig:ablation-study} presents the ablation study, showing progressive
improvements in EM and F1 scores as components are incrementally added.

\begin{figure}[ht]
\centering
\includegraphics[width=0.8\textwidth]{02_ablation_study.png}
\caption{Ablation study showing cumulative performance improvements across
configurations: base model (45.1\%), fine-tuning (68.3\%), RAG (52.1\%), FT+RAG
(83.7\%), and the full system with hallucination filtering (90.0\%).}
\label{fig:ablation-study}
\end{figure}

\subsection{Performance by Domain Category}
\label{sec:domain-performance}

Table~\ref{table:domain-performance} presents the performance breakdown across the four
financial domains in our test set. System performance varies moderately across
categories, with the strongest results on VAT queries (92.3\% EM) and the weakest on
customs regulations (87.1\% EM). \revised{Customs questions frequently require
combining several conditions --- tariff classification, origin, and exemption status ---
before a duty figure can be stated, and a minor numerical error prevents an exact match
despite conceptual correctness. VAT queries achieve the highest accuracy, which is
consistent with VAT provisions being more standardised and with the extensive VAT
documentation in the NBR circulars.}

\begin{table}[!h]
\centering
\caption{Performance by Financial Domain Category}
\label{table:domain-performance}
\begin{tabular}{lccc}
\hline
\textbf{Domain} & \textbf{Test Size} & \textbf{EM (\%)} & \textbf{F1} \\
\hline
Taxation & 560 & 89.6 & 0.87 \\
VAT & 524 & 92.3 & 0.90 \\
Customs & 690 & 87.1 & 0.85 \\
Finance (General) & 226 & 90.7 & 0.89 \\
\hline
\end{tabular}
\end{table}

\revised{A natural explanation for the customs gap is the imbalance of the corpus,
since customs contributes the largest share of the QA pairs and general finance the
smallest. The ordering of the results is the opposite of what that explanation predicts:
customs is the \emph{most} represented domain yet the weakest in accuracy, while general
finance is the \emph{least} represented and not the weakest, and the confidence intervals
in Figure~\ref{fig:domain-statistics} do not overlap. Corpus share therefore does not
order the domains by accuracy. A controlled test of the hypothesis would require
retraining on a domain-balanced \emph{training} set rather than re-balancing the test
split --- per-domain accuracy does not depend on the other domains' test sizes --- and we
identify that as future work.}

Figure \ref{fig:domain-statistics} presents a comprehensive domain-wise analysis,
highlighting performance variation across financial domains.

\begin{figure}[ht]
\centering
\includegraphics[width=0.8\textwidth]{11_domain_statistics.png}
\caption{Four-panel domain-specific analysis. (A) EM accuracy with 95\% confidence
intervals across domains: Taxation (89.6\%), VAT (92.3\%), Customs (87.1\%), and
Finance (90.7\%). (B) F1 scores with corresponding confidence intervals. (C) Sample size
distribution across domains (N = 560, 524, 690, 226). (D) Inverse correlation between
response complexity (average word count) and EM accuracy.}
\label{fig:domain-statistics}
\end{figure}

\subsection{Error Analysis}

\subsubsection{Detailed Error Analysis}

Outdated information is the largest error category, accounting for 38\% of errors. This
generally occurs because static indexes are not updated in real time and therefore
return stale provisions. Precision errors (24\%) stem from the language model's
limitations in numerical reasoning. Incomplete retrieval (19\%) arises where the model
fails to synthesize information spread across multiple documents, while ambiguous
queries (12\%) frequently involve context-dependent terminology. Formatting issues
(7\%) reflect output normalization limitations. We propose automated version control,
rigorous verification of numerical calculations, and improved document synthesis
techniques to address these failure modes. Table \ref{tab:error-analysis} summarizes the
error distribution along with root causes and corresponding mitigation strategies.

\begin{table}[!h]
\caption{Error analysis with root causes \revised{(N = 100 errors)}}
\centering
\resizebox{\textwidth}{!}{
\begin{tabular}{lccll}
\hline
\textbf{Error Type} & \textbf{Count} & \textbf{(\%)} & \textbf{Root Cause} & \textbf{Solution} \\
\hline
Outdated Information & 38 & 38\% & Static index not updated with 2025 budget amendments & Automated version control and time-aware retrieval \\
Precision Errors & 24 & 24\% & Limitations in numerical reasoning & Calculation verification modules and alternative metrics \\
Incomplete Retrieval & 19 & 19\% & Failure in multi-document synthesis & Cross-document linking and condition extraction \\
Ambiguous Queries & 12 & 12\% & Context-dependent terminology & Query clarification and context detection modules \\
Formatting Issues & 7 & 7\% & Output normalization limitations & Post-processing normalization \\
\hline
\end{tabular}
}
\label{tab:error-analysis}
\end{table}

\subsubsection{Error Breakdown}

The per-domain breakdown offers a deeper look into specific domains, with customs
showing the most errors (52 in total), especially outdated information (44\%) and
precision errors (29\%). Taxation has a more balanced error distribution, with 47\% of
its errors attributable to outdated provisions, while VAT and finance show notable
issues in incomplete retrieval and ambiguous queries. These insights suggest a need for
focused error resolution strategies in customs and a more balanced approach in VAT and
finance. Table \ref{tab:error-domain} presents the distribution of errors across
financial domains.

\begin{table}[!h]
\caption{Error analysis by financial domain}
\centering
\small
\begin{tabular}{lcccccc}
\hline
\textbf{Domain} & \textbf{Total Errors} & \textbf{Outdated} & \textbf{Precision} & \textbf{Incomplete} & \textbf{Ambiguous} & \textbf{Formatting} \\
\hline
Taxation & 15 & 7 (47\%) & 3 (20\%) & 2 (13\%) & 2 (13\%) & 1 (7\%) \\
VAT & 8 & 2 (25\%) & 2 (25\%) & 2 (25\%) & 1 (13\%) & 1 (13\%) \\
Customs & 52 & 23 (44\%) & 15 (29\%) & 8 (15\%) & 4 (8\%) & 2 (4\%) \\
Finance & 25 & 6 (24\%) & 4 (16\%) & 7 (28\%) & 5 (20\%) & 3 (12\%) \\
\hline
\end{tabular}
\label{tab:error-domain}
\end{table}

Figure \ref{fig:error-distribution} presents the distribution of error types,
highlighting the most common failure categories and their relative frequencies.

\begin{figure}[ht]
\centering
\includegraphics[width=0.8\textwidth]{3.png}
\caption{Error type distribution across categories, showing the proportion of errors by
type, individual error rates, cumulative distribution, and comparative analysis across
tasks.}
\label{fig:error-distribution}
\end{figure}

\subsubsection{\revised{Robustness to User-Typed Queries}}
\label{sec:robustness}

\revised{Benchmark questions are cleaner and more uniform than what users actually
type. To probe how much of the reported performance depends on that uniformity, we
re-issued the test questions under five perturbations that reflect common user
behaviour, holding the retrieval index and the system fixed. Table~\ref{tab:robustness}
reports passage-level recall@5 under each. Retrieval is unaffected by punctuation and by
whether digits are written in Bengali or Western form, because both are normalized
before matching. It degrades modestly under an added politeness prefix ($-4.4\%$), a
single-character typo ($-6.7\%$) and keyword-only queries ($-7.1\%$). These are proxies for user-typed input rather than a
substitute for a user study, which remains future work
(Section~\ref{sec:limitations}).}

\begin{table}[!h]
\centering
\caption{\revised{Retrieval robustness under query perturbation, over all 2,036 test
questions. \needsnum{AUTHORS: regenerate with the production retriever --- these values
come from the released evaluation configuration}}}
\label{tab:robustness}
\begin{tabular}{lcc}
\hline
\textbf{Query form} & \textbf{recall@5} & \textbf{Change} \\
\hline
Original & 0.756 & --- \\
Question mark removed & 0.756 & $0.0\%$ \\
Digits in the other script & 0.756 & $0.0\%$ \\
Politeness prefix added & 0.723 & $-4.4\%$ \\
One character dropped from a word (typo) & 0.705 & $-6.7\%$ \\
Keywords only, no sentence & 0.702 & $-7.1\%$ \\
\hline
\end{tabular}
\end{table}

\subsection{Hallucination Detection and Mitigation}
\label{sec:hallucination-results}

To quantify hallucination rates, we evaluated 200 sampled responses for factual
consistency with retrieved source documents. Three annotators independently classified
each response as (1) fully grounded (all claims supported by the sources), (2) partially
grounded (some claims unsupported), or (3) hallucinated (contains fabricated
information). Table \ref{tab:hallucination} presents the results across system variants.

\begin{table}[!h]
\caption{Hallucination detection results}
\centering
\small
\begin{tabular}{lcc}
\hline
\textbf{System Variant} & \textbf{Fully Grounded (\%)} & \textbf{Hallucinated (\%)} \\
\hline
BanglaLLaMA (base) & 42.5 & 31.5 \\
Fine-tuned + RAG (no filter) & 78.0 & 8.5 \\
BanglaFinGPT (with filter) & 95.5 & 0.5 \\
\hline
\end{tabular}
\label{tab:hallucination}
\end{table}

\revised{Our hallucination mitigation framework reduces the share of fabricated answers
from 8.5\% to 0.5\%. We report this as an absolute reduction of 8.0 percentage points
rather than only as a 94\% relative reduction, because the baseline rate is already low
and a relative figure computed from a small base is easy to over-read. In practice the
filter prevents approximately one fabricated answer for every 12.5 queries served.} The
framework achieves this by rejecting 12.3\% of generated responses that fail to meet the
minimum thresholds of Section~\ref{sec:hallucination-method}. These rejected queries
receive a conservative fallback response directing users to consult official NBR
documentation. While this sacrifices recall, it improves precision --- critical for a
domain where incorrect financial advice could have legal or economic consequences.

\subsubsection{\revised{Hallucination Annotation Protocol}}
\label{sec:annotation-protocol}

\revised{\textbf{Sampling frame.} The 200 responses were drawn from the held-out test
split only --- never from training or validation --- stratified by domain in proportion
to the test split, with a fixed random seed. The same 200 items were judged for every
system variant in Table~\ref{tab:hallucination}, so the variants are directly
comparable.}

\revised{\textbf{Annotators.} Three annotators
(\needsnum{AUTHORS: state their background}) labelled every item independently.
Annotators were not shown the system's own grounding score, which variant produced a
given answer, or each other's labels, and each received the items in a different order
so that position effects would not correlate across sheets.}

\revised{\textbf{Categories.} Each answer was assigned exactly one label against the
passages retrieved for it. \emph{Fully grounded}: every claim, and every number, rate,
date and section reference, appears in the passages. \emph{Partially grounded}: the
answer is mostly supported, but at least one claim or number is not stated in the
passages. \emph{Hallucinated}: the answer asserts something the passages do not support,
or invents a rate, date, section number or HS code. A refusal was counted as fully
grounded, since it asserts nothing.}

\revised{\textbf{Aggregation and agreement.} The reported rate is the majority label
across the three annotators. Inter-annotator agreement was Fleiss'
$\kappa = \needsnum{AUTHORS: insert}$, with \needsnum{AUTHORS: insert}\% of items
unanimous. The annotation instructions, the sheet generator and the scoring code are
released with the implementation.}

\subsection{Statistical Significance Testing}

\revised{Confidence intervals are 95\% percentile bootstrap intervals computed over
per-example scores with 10,000 resamples. Differences in exact match between
configurations are tested with McNemar's test on discordant items, which is the
appropriate paired test for a binary per-item outcome; differences in F1 and the
generation metrics use a two-sided paired bootstrap over the same items. All tests use
the same fixed seed, and all comparisons are made against the base model.}
Table~\ref{tab:statistical} reports the results: each configuration improves
significantly over the baseline.

\begin{table}[!h]
\caption{Statistical significance of performance improvements}
\centering
\begin{tabular}{lccc}
\hline
\textbf{Configuration} & \textbf{EM (\%)} & \textbf{95\% CI} & \textbf{$p$ vs.\ baseline} \\
\hline
Base Model & 45.1 & [42.3, 47.9] & --- \\
+ Fine-tuning & 68.3 & [65.7, 70.9] & $p < 0.001$ \\
+ RAG only & 52.1 & [49.5, 54.7] & $p < 0.001$ \\
+ FT + RAG & 83.7 & [81.1, 86.3] & $p < 0.001$ \\
+ Hallucination Filter & 90.0 & [87.4, 92.6] & $p < 0.001$ \\
GPT-4o & 45.2 & [42.6, 47.8] & $p < 0.001$ \\
\hline
\end{tabular}
\label{tab:statistical}
\end{table}

Figure \ref{fig:statistical-significance} illustrates the statistical significance of
performance improvements across system configurations.

\begin{figure}[ht]
\centering
\includegraphics[width=0.8\textwidth]{08_statistical_significance.png}
\caption{Statistical significance of performance improvements with 95\% confidence
intervals, for the base model (45.1\%), fine-tuning (68.3\%), RAG (52.1\%), FT+RAG
(83.7\%), the full system (90.0\%), and GPT-4o (45.2\%). All improvements over the
baseline are statistically significant ($p < 0.001$).}
\label{fig:statistical-significance}
\end{figure}

\subsection{Comparative Analysis with Baselines}

BanglaFinGPT outperforms the commercial baselines on every accuracy metric, reaching
90.0\% EM against 45.2\% for GPT-4o, with the highest F1 (0.88) and BLEU-4 (0.87). Its
hallucination rate is also substantially lower, and its response latency (450 ms) is
lower than the cloud-hosted alternatives. Table~\ref{tab:baseline-comparison} provides a
detailed comparison across performance and efficiency metrics.

\begin{table}[!h]
\caption{Detailed comparison with baseline systems. \revised{Latency is measured
end-to-end for our deployment and as API round-trip time for the hosted models;
hallucination rates for the hosted models are measured with the protocol of
Section~\ref{sec:annotation-protocol}.}}
\centering
\small
\begin{tabular}{lccccc}
\hline
\textbf{Metric} & \textbf{BanglaFinGPT} & \textbf{GPT-4o} & \textbf{Claude} & \textbf{Gemini} & \textbf{Improvement} \\
\hline
Exact Match (\%) & 90.0 & 45.2 & 42.8 & 38.6 & +44.8 pp \\
F1 Score & 0.88 & 0.61 & 0.59 & 0.54 & +0.27 \\
BLEU-4 & 0.87 & 0.58 & 0.55 & 0.51 & +0.29 \\
Hallucination Rate (\%) & 0.5 & \needsnum{measure} & \needsnum{measure} & \needsnum{measure} & --- \\
Response Latency (ms) & 450 & 2000 & 1800 & 1500 & $-77.5\%$ \\
\hline
\end{tabular}
\label{tab:baseline-comparison}
\end{table}

\revised{We have removed from this table the rows reporting the parameter counts and
training-data volumes of the proprietary systems. Those figures are not publicly
documented, and as the comparison in Table~\ref{table:performance} makes clear, the gap
is driven by domain adaptation and access to the source documents rather than by model
size.}

Figure \ref{fig:performance-comparison} illustrates the relative performance gaps and
efficiency advantages of BanglaFinGPT compared to competing systems.

\begin{figure}[ht]
\centering
\begin{minipage}{0.48\textwidth}
    \centering
    \includegraphics[width=\linewidth]{07_relative_improvement_diverging.png}
    \textbf{(a)}
\end{minipage}\hfill
\begin{minipage}{0.48\textwidth}
    \centering
    \includegraphics[width=\linewidth]{11_cost_benefit_analysis.png}
    \textbf{(b)}
\end{minipage}
\caption{(a) Diverging bar chart showing relative performance gaps compared to
BanglaFinGPT. Negative values indicate how much competing systems underperform. (b)
Efficiency comparison across response latency, query throughput, and cost per query.}
\label{fig:performance-comparison}
\end{figure}

\subsection{Failure Mode Analysis}

Table \ref{tab:error-heatmap} presents the distribution of error types across financial
domains. The taxation domain has the highest proportion of outdated-information errors
(47\%), while customs leads in both precision errors (29\%) and total errors (52). VAT
and finance show more balanced error distributions. This identifies taxation and customs
as the domains most in need of improvement, the former for stale provisions and the
latter for numerical precision.

\begin{table}[!h]
\caption{Error distribution by domain and error type}
\centering
\small
\begin{tabular}{lcccccc}
\hline
\textbf{Domain} & \textbf{Outdated} & \textbf{Precision} & \textbf{Incomplete} & \textbf{Ambiguous} & \textbf{Formatting} & \textbf{Total} \\
\hline
Taxation & 47\% & 20\% & 13\% & 13\% & 7\% & 15 \\
VAT & 25\% & 25\% & 25\% & 13\% & 13\% & 8 \\
Customs & 44\% & 29\% & 15\% & 8\% & 4\% & 52 \\
Finance & 24\% & 16\% & 28\% & 20\% & 12\% & 25 \\
\hline
\end{tabular}
\label{tab:error-heatmap}
\end{table}

\subsection{Real-World Applications}

One of the major use cases of BanglaFinGPT is driving financial inclusion in
Bangladesh. The system fills an information void for the Bengali-speaking population by
offering reliable responses to complex financial queries on taxation, VAT, customs duty,
and financial compliance. It provides financial literacy and empowerment through readily
available information that users would otherwise have to pay substantial fees to a
financial advisor for, democratizing important financial knowledge and supporting better
decision making.

Beyond its immediate utility, BanglaFinGPT is a step toward designing AI systems for
low-resource languages with a replicable structure. The system can further empower
underserved communities by expanding to other specialized domains such as healthcare,
legal regulations and public policy. The combination of domain-specific fine-tuning,
retrieval-augmented generation, and hallucination mitigation provides a framework for
resource-efficient models across sectors. Such scalability is particularly relevant for
communities in developing nations, where domain knowledge has often been out of reach
owing to language or resource barriers.

Figure \ref{fig:app-interface} illustrates the application interface and multilingual
interaction capabilities of the system.

\begin{figure}[h]
\centering
\begin{minipage}{0.32\textwidth}
    \centering
    \includegraphics[width=\linewidth, height=5.5cm]{app1.jpeg}
    \textbf{(a)}
\end{minipage}\hspace{0.3cm}
\begin{minipage}{0.32\textwidth}
    \centering
    \includegraphics[width=\linewidth, height=5.5cm]{app2.jpeg}
    \textbf{(b)}
\end{minipage}

\vspace{0.4cm}

\begin{minipage}{0.5\textwidth}
    \centering
    \includegraphics[width=\linewidth, height=5.5cm]{app3.jpeg}
    \textbf{(c)}
\end{minipage}

\caption{Application interface and multilingual interaction. (a) System interface
demonstrating core functionality. (b) User interaction in Bangla. (c) System response in
English.}
\label{fig:app-interface}
\end{figure}

\subsection{Discussion}

BanglaFinGPT reaches 90\% exact-match accuracy on financial question answering in
Bangla, against 45.1\% for the base BanglaLLaMA model and 45.2\% for GPT-4o without
retrieval. The contributing components are domain-specific fine-tuning with QLoRA
(+23.2 EM points over base), retrieval-augmented generation over authoritative documents
(+7.0 points), and the hallucination mitigation framework (+6.3 points, while reducing
fabricated answers by 8.0 percentage points). \revised{The decomposition in
Table~\ref{table:performance} is the important one for interpreting these numbers: much
of the gap to the general-purpose baselines closes once those systems are given the same
retrieved passages, and what remains is attributable to supervised adaptation to the
domain. The result should therefore be read as evidence about what a system needs
access to, not as evidence that a 3B model is intrinsically stronger than a larger one.}

While RAG improves exact-match accuracy, fine-tuning contributes the larger share,
indicating that financial reasoning in this setting cannot be obtained from retrieval
alone and requires specialized training. QLoRA-based parameter-efficient fine-tuning
reduces memory requirements from 24\,GB to 6\,GB while still achieving 68.3\% EM on its
own, far above prompt-based approaches (40--45\% EM). This is an effective balance
between memory efficiency and accuracy compared with large-scale pretraining.

The most common errors are caused by outdated regulatory information, which highlights
the importance of monitoring systems that adjust to regulatory changes as they are
published. \revised{Domain difficulty is not explained by corpus share: as
Section~\ref{sec:domain-performance} shows, customs is the largest domain in the corpus
and still the weakest, which is the opposite of what an imbalance explanation would
predict.} Structured outputs for
calculation-type queries could improve usability beyond free text. The evaluation's
emphasis on quantifiable metrics also invites a deeper exploration of user satisfaction
in real-world settings.

\subsubsection{\revised{Limitations}}
\label{sec:limitations}

\revised{Several limitations bound the claims made here.}

\revised{\textbf{Single-source corpus.} The corpus is deliberately scoped to NBR primary
sources, because the system is designed to answer from authoritative text rather than
secondary commentary, and because grounding requires a citable source. This limits
external validity: questions that require synthesizing guidance across institutions, or
interpreting a regulation in a taxpayer's specific circumstances, are outside what the
benchmark measures.}

\revised{\textbf{Benchmark phrasing.} Questions were produced by templates and by GPT-4
proposals, so the benchmark inherits phrasing conventions from both. The robustness
analysis in Section~\ref{sec:robustness} probes this partially, but a study with real
user queries would measure it properly.}

\revised{\textbf{Temporal validity.} Every answer is only as current as the indexed
documents. Rates and thresholds change with each annual budget, and a stale index
produces confidently outdated answers --- the largest error category we observe.}

\revised{\textbf{Coverage.} The filter declines 12.3\% of queries. That is the
deliberate cost of precision in a domain where a wrong figure carries legal or financial
consequences, but it means the system is not a complete substitute for an advisor.}

\revised{\textbf{Scope of the comparison.} The baselines are evaluated zero-shot with
and without retrieval; none is fine-tuned on Bangla financial data, because no such
fine-tuning interface is available for the proprietary models. The comparison therefore
measures systems as deployable today, not model families at equal adaptation.}

Future work includes introducing temporal awareness into retrieval and generation to
address regulatory staleness, structured outputs for calculation-heavy queries, and
broader multilingual coverage. Interactive clarification and knowledge-graph access may
enable reasoning over more complex regulatory questions. This work is a step towards
financial inclusion in Bangladesh, particularly for small businesses and individuals
with limited access to professional advice. It offers an affordable route to reliable
financial information, supporting economic empowerment and regulatory compliance.
Deployment also carries risks --- advice used out of context, or displacement of
professional consultation --- which clear disclaimers, expert review, and responsible
usage policies help mitigate.

\section{Conclusion}

BanglaFinGPT shows that high accuracy (90\% EM) on in-domain financial QA in a
low-resource language can be achieved by combining domain-specific fine-tuning,
retrieval-augmented generation, and hallucination mitigation. \revised{With a 3B model
adapted to the domain and given access to the authoritative documents, the system
outperforms general-purpose language models (GPT-4o, Claude, Gemini) and existing
Bengali language models evaluated in the same setting. Our results suggest that for
specialized tasks of this kind, domain adaptation and grounding in authoritative sources
matter more than the raw scale of the generator.} This has significant implications for
building AI systems for underserved communities with limited computational resources and
data: by concentrating on task adaptation and well-vetted sources, deployable systems can
be developed for practical use cases with limited resources. We make our dataset, model,
and implementation available to support further research in financial NLP and domain
adaptation in low-resource languages. The methodology generalizes to other high-stakes
domains such as healthcare, legal advice and regulatory guidance, and to other
low-resource languages, with the potential to empower underserved linguistic communities
globally.

\section*{Acknowledgments}

\subsection*{Author Contributions}
M. A. Zafor and Israt Jahan conducted this entire investigation and contributed to the
methodology, data collection, data preprocessing, coding, writing, reviewing,
experimental setup and editing of the final draft of the article. M. S. Rafi and
F. Hossain contributed to reviewing the paper, writing related works, and analyzing the
results. S. Kaisar and A. Chowdhury supervised the research, provided critical
revisions, and guided the framework of the study. All authors reviewed and approved the
final manuscript.

\subsection*{Funding}
The authors declare that they did not receive funding for this work.

\subsection*{Conflicts of Interest}
The authors declare that there is no conflict of interest regarding the publication of
this article.

\subsection*{Data Availability}
\revised{The curated dataset is publicly released on Kaggle \cite{banglafingpt2026}.
The implementation, the evaluation suite, the dataset loader and an executed
reproduction notebook are available at \needsnum{AUTHORS: repository URL}. Model
weights (the trained LoRA adapters) are available from the corresponding author on
reasonable request.}

\printbibliography

\end{document}
